\section{Introduction}
Distributed systems replicate state across machines to provide fault
tolerance, and consensus protocols coordinate replicas so that their
states remain consistent. Two design families provide contrasting
approaches to this coordination. Leader-based protocols such as
Raft~\cite{ongaro2014raft} direct client operations through an elected
leader, which orders log entries and coordinates replication; leader
failure can therefore require an election before normal processing
resumes. Leaderless protocols such as EPaxos~\cite{moraru2013epaxos}
allow replicas to propose operations independently and determine
dependencies among them, avoiding a single ordering point at the cost of
dependency tracking and conflict handling.

These architectural differences motivate empirical comparison, but a
simple protocol-level ranking is difficult to interpret. Measured
performance depends on replica count, client concurrency, workload
composition, conflict rate, and communication conditions. It is also
affected by implementation mechanisms such as persistence, scheduling,
transport, and message handling. A measured difference between two
implementations therefore cannot automatically be attributed to the
consensus algorithm alone. This is particularly important when comparing
leader-based and leaderless designs: the presence of a leader is a
structural difference, but it does not by itself determine the throughput,
latency, or failure behaviour of a concrete implementation.

In this paper, we experimentally compare Raft and EPaxos as
representative leader-based and leaderless consensus protocols. We
evaluate throughput, latency, resource utilisation, and failure behaviour
under controlled variations in replica count, client concurrency, conflict
rate, read/write composition, and communication conditions. We further
evaluate multiple implementations and matched persistence configurations
to examine whether observed differences remain consistent across
implementation choices.

We focus on the following research questions:
\begin{itemize}
\item \textbf{RQ1:} How do throughput and latency change with workload
and system conditions, including read/write ratio, client concurrency,
conflict rate, and replica count?
\item \textbf{RQ2:} How do Raft and EPaxos differ in failure behaviour
and in their sensitivity to communication cost?
\item \textbf{RQ3:} To what extent are the observed performance
differences between Raft and EPaxos attributable to implementation choices
rather than to the protocol families themselves?
\end{itemize}

We make the following contributions:
\begin{itemize}
\item \textbf{Reproducible benchmark infrastructure:} the Consensus Lab
executes Raft and EPaxos in a common containerised environment through a
common client harness and benchmark-facing interface.
\item \textbf{Sensitivity-oriented evaluation:} throughput and latency
are measured across controlled variations in replica count, client
concurrency, conflict rate, read/write composition, and communication
conditions. This allows the evaluation to examine changes across
conditions rather than relying on a single benchmark configuration.
\item \textbf{Failure and communication analysis:} we characterise
availability under Raft leader failure, Raft follower failure, and EPaxos
replica failure, and evaluate the effect of increasing communication cost
using controlled network-delay experiments.
\item \textbf{Implementation-level analysis:} multiple implementations
of the two protocol families are evaluated, including matched persistence
configurations and communication-cost experiments. The results show that
implementation mechanisms can materially change the measured relationship
between Raft and EPaxos.
\item \textbf{Open infrastructure:} the lab, configurations, and analysis
pipeline are released at
\url{https://github.com/Sephy314/Leader-consensus-comparison}.
\end{itemize}

The experiments run in a local containerised environment, with controlled
communication delay emulated at the operating-system network layer.
Therefore, the results are measurements of the evaluated implementations
under the specified conditions rather than universal production-
performance estimates. In particular, the communication experiments
characterise sensitivity to increased communication cost but do not
reproduce all properties of a real wide-area deployment, such as
heterogeneous routing, packet loss, or independent network paths.

The remainder of this paper is organised as follows.
Section~\ref{sec:background} reviews leader-based and leaderless consensus
and related comparative studies. Section~\ref{sec:setup} describes the
Consensus Lab, the implementations, and the experiment suite.
Section~\ref{sec:evaluation} presents the results.
Section~\ref{sec:discussion} interprets them with respect to the research
questions. Section~\ref{sec:threats} discusses threats to validity, and
Section~\ref{sec:conclusion} summarises the findings.
